% --- IMPORTS ---
\documentclass[leqno]{article}
\usepackage{graphicx}
\usepackage{dsfont}
\usepackage{mathtools}
\usepackage{amssymb}
\usepackage{amsmath}
\usepackage{amsthm}
\usepackage{float}
\usepackage{ragged2e}
\usepackage{tensor}
\usepackage{fancyhdr}
\usepackage{empheq}
\usepackage[most]{tcolorbox}
\usepackage{enumitem}
\usepackage{dirtytalk}
\usepackage{marginnote}
\usepackage{microtype}
\usepackage[
  a4paper,
  left=0.8in,
  right=1in,
  top=1in,
  bottom=0.5in,
  headheight=14pt,
  headsep=0.4in
]{geometry}
\setlength{\parindent}{0cm}
% --- TikZ Packages ---
\usepackage{pgfplots}
\pgfplotsset{compat=1.18}
\usepgfplotslibrary{fillbetween, polar}
\usetikzlibrary{calc, positioning}
\usetikzlibrary{angles, quotes}
\usetikzlibrary{arrows.meta}
\usetikzlibrary{decorations.pathreplacing, decorations.markings}
\usetikzlibrary{patterns}
\usetikzlibrary{intersections}
% --- URL Packages ---
\usepackage[colorlinks=true,linkcolor=red,citecolor=magenta,urlcolor=black]{hyperref}%
\urlstyle{same}
% --- END OF IMPORTS ---

% --- CUSTOM COMMANDS ---
\RenewDocumentCommand{\marks}{o m}{%
  \IfValueTF{#1}%
    {\marginnote{\raggedright\textbf{[#2]}}[#1]}%
    {\marginnote{\raggedright\textbf{[#2]}}}%
}
\newcommand{\vspacerule}[2]{%
  \par\vspace*{#1}%
  \noindent\hrulefill\par
  \vspace*{#2}%
}
\definecolor{scarlet}{RGB}{205, 40, 75}
\newcommand{\questionitem}{%
  \item\phantomsection
  \addcontentsline{toc}{section}{Question \number\value{enumi}}%
}
\renewcommand{\contentsname}{Questions}
% --- END OF CUSTOM COMMANDS ---

% --- HEADER CONFIG ---
\pagestyle{fancy}
\fancyhead{}
\fancyfoot{}
\renewcommand{\headrulewidth}{0.9pt}
\lhead{\textit{Solutions} - Maclaurin Series}
\chead{}
\rhead{\href{https://danieljsmith.org}{danieljsmith.org}}
% --- END OF HEADER CONFIG ---

\begin{document}
\vspace*{20pt}
\tableofcontents
\newpage
\begin{enumerate}[
  label=\textbf{\arabic*.},
  leftmargin=*,
  topsep=0pt,
  partopsep=0pt,
  parsep=0pt
]

% ---- Question 1 ----
\questionitem
% [Source: _QBT___Solns__Maclaurin_Series, Q1]

\begin{enumerate}[label=\textbf{(\alph*)}]
  \item
  \begin{enumerate}[label=(\roman*)]
    \item Given that $f(x)=e^{-2x}$, find $f'(x)$ and $f''(x)$. \marks{2}
    \item Hence, find the first three terms of the Maclaurin series for $e^{-2x}$. \marks{2} \\
  \end{enumerate}
  \item Hence, using a suitable value of $x$, show that \[e^{-1/4}\approx \frac{25}{32}\]. \marks{2} \\
\end{enumerate}

\vspace*{10pt}

\begin{tcolorbox}[colback=white, colframe=scarlet, title={\textbf{Solution}}, breakable, grow to left by=6mm, grow to right by=10mm]
\begin{enumerate}[label=\textbf{(\alph*)}]
  \item
  \begin{enumerate}[label=(\roman*)]
    \item Differentiate \(f(x)=e^{-2x}\) using the chain rule:
    \[
    f'(x)=(-2)e^{-2x}=-2e^{-2x}
    \]
    Differentiate again:
    \[
    f''(x)=(-2)(-2)e^{-2x}=4e^{-2x}
    \]
    So
    \[
    f'(x)=-2e^{-2x}, \qquad f''(x)=4e^{-2x}
    \]

    \item The Maclaurin series is
    \[
    f(x)=f(0)+xf'(0)+\frac{x^2}{2!}f''(0)+\cdots
    \]
    Now
    \[
    f(0)=e^0=1,\qquad f'(0)=-2e^0=-2,\qquad f''(0)=4e^0=4
    \]
    So
    \begin{align*}
    e^{-2x}
    &=1+x(-2)+\frac{x^2}{2}(4)+\cdots \\
    &=1-2x+2x^2+\cdots
    \end{align*}
    Hence the first three terms are
    \[
    1-2x+2x^2
    \]
  \end{enumerate}

  \item To get \(e^{-1/4}\), choose \(x\) so that
  \[
  -2x=-\frac14
  \]
  Hence
  \[
  x=\frac18
  \]
  Using the Maclaurin expansion found in part (a),
  \[
  e^{-2x}\approx 1-2x+2x^2
  \]
  Substituting \(x=\frac18\),
  \begin{align*}
  e^{-1/4}
  &=e^{-2(1/8)} \\
  &\approx 1-2\left(\frac18\right)+2\left(\frac18\right)^2 \\
  &=1-\frac14+2\cdot\frac{1}{64} \\
  &=1-\frac14+\frac{1}{32} \\
  &=\frac{25}{32}
  \end{align*}
  So
  \[
  e^{-1/4}\approx \frac{25}{32}
  \]
\end{enumerate}
\end{tcolorbox}


\newpage

% ---- Question 2 ----
\questionitem
% [Source: _QBT___Solns__Maclaurin_Series, Q2]

\begin{enumerate}[label=(\roman*)]
  \item Use the Maclaurin series for $\sin x$ and $\cos x$ to work out the series expansion of \[\sin 2x(\cos x-\cos 3x)\] up to and including the term in $x^5$. \marks{4} \\
  \item Hence show that, in exact surd form, an approximation to the least positive root of the equation \[2\sin 2x(\cos x-\cos 3x)=x\] is \[\sqrt{\frac{4-\sqrt{10}}{12}}
\]\marks[-20pt]{3} \\
\end{enumerate}

\vspace*{10pt}

\begin{tcolorbox}[colback=white, colframe=scarlet, title={\textbf{Solution}}, breakable, grow to left by=6mm, grow to right by=10mm]
\begin{enumerate}[label=\textbf{(\roman*)}]
\item Using the standard Maclaurin series,
\[
\sin x=x-\frac{x^3}{3!}+\frac{x^5}{5!}+\cdots,
\qquad
\cos x=1-\frac{x^2}{2!}+\frac{x^4}{4!}+\cdots
\]

First expand \(\sin 2x\):
\begin{align*}
\sin 2x
&=2x-\frac{(2x)^3}{3!}+\frac{(2x)^5}{5!}+\cdots\\
&=2x-\frac{8x^3}{6}+\frac{32x^5}{120}+\cdots\\
&=2x-\frac{4}{3}x^3+\frac{4}{15}x^5+\cdots
\end{align*}

Now expand \(\cos x-\cos 3x\):
\begin{align*}
\cos x&=1-\frac{x^2}{2}+\frac{x^4}{24}+\cdots\\
\cos 3x&=1-\frac{(3x)^2}{2}+\frac{(3x)^4}{24}+\cdots\\
&=1-\frac{9x^2}{2}+\frac{81x^4}{24}+\cdots
\end{align*}

So
\begin{align*}
\cos x-\cos 3x
&=\left(1-\frac{x^2}{2}+\frac{x^4}{24}\right)-\left(1-\frac{9x^2}{2}+\frac{81x^4}{24}\right)+\cdots\\
&=4x^2-\frac{80x^4}{24}+\cdots\\
&=4x^2-\frac{10}{3}x^4+\cdots
\end{align*}

Now multiply:
\begin{align*}
\sin 2x(\cos x-\cos 3x)
&=\left(2x-\frac{4}{3}x^3+\frac{4}{15}x^5+\cdots\right)\left(4x^2-\frac{10}{3}x^4+\cdots\right)
\end{align*}

Keeping terms up to \(x^5\),
\begin{align*}
\sin 2x(\cos x-\cos 3x)
&=(2x)(4x^2)+(2x)\left(-\frac{10}{3}x^4\right)+\left(-\frac{4}{3}x^3\right)(4x^2)+\cdots\\
&=8x^3-\frac{20}{3}x^5-\frac{16}{3}x^5+\cdots\\
&=8x^3-12x^5+\cdots
\end{align*}

Hence,
\[
\sin 2x(\cos x-\cos 3x)=8x^3-12x^5+\cdots
\]

\item From part (i),
\[
2\sin 2x(\cos x-\cos 3x)\approx 2(8x^3-12x^5)=16x^3-24x^5
\]

So the equation
\[
2\sin 2x(\cos x-\cos 3x)=x
\]
becomes approximately
\[
16x^3-24x^5=x
\]

Rearranging,
\begin{align*}
16x^3-24x^5-x&=0\\
x(16x^2-24x^4-1)&=0\\
x(24x^4-16x^2+1)&=0
\end{align*}

Since we want the least positive root, \(x\neq 0\), so
\[
24x^4-16x^2+1=0
\]

Let \(y=x^2\). Then
\[
24y^2-16y+1=0
\]

Using the quadratic formula,
\begin{align*}
y&=\frac{16\pm\sqrt{(-16)^2-4(24)(1)}}{2(24)}\\
&=\frac{16\pm\sqrt{256-96}}{48}\\
&=\frac{16\pm\sqrt{160}}{48}\\
&=\frac{16\pm 4\sqrt{10}}{48}\\
&=\frac{4\pm\sqrt{10}}{12}
\end{align*}

For the least positive root, take the smaller value:
\[
x^2\approx \frac{4-\sqrt{10}}{12}
\]

Hence
\[
x\approx \sqrt{\frac{4-\sqrt{10}}{12}}
\]

So the least positive root is approximately
\[
\sqrt{\frac{4-\sqrt{10}}{12}}\approx 0.264
\]
\end{enumerate}
\end{tcolorbox}


\newpage

% ---- Question 3 ----
\questionitem
% [Source: _QBT___Solns__Maclaurin_Series, Q3]

\begin{enumerate}[label=\textbf{(\alph*)}]
  \item Starting from the series given in the formulae booklet, show that the general term of the Maclaurin series for
  \[
  2(1-\cos x)-x\sin x
  \]
  is
  \[
  (-1)^r\frac{2(r-1)}{(2r)!}x^{2r}
  \]
  \marks[-20pt]{4}
  \item Show that
  \[
  \lim_{x\to 0}\left[\frac{2(1-\cos x)-x\sin x}{(1-\cos x)^2}\right]=\frac{1}{3}
  \]
  \marks[-30pt]{4}
\end{enumerate}

\vspace*{10pt}

\begin{tcolorbox}[colback=white, colframe=scarlet, title={\textbf{Solution}}, breakable, grow to left by=6mm, grow to right by=10mm]
\begin{enumerate}[label=\textbf{(\alph*)}]
  \item From the formula booklet,
\[
\cos x=\sum_{r=0}^{\infty}(-1)^r\frac{x^{2r}}{(2r)!},
\qquad
\sin x=\sum_{r=0}^{\infty}(-1)^r\frac{x^{2r+1}}{(2r+1)!}
\]

So
\begin{align*}
2(1-\cos x)
&=2\left(1-\sum_{r=0}^{\infty}(-1)^r\frac{x^{2r}}{(2r)!}\right)\\
&=2\sum_{r=1}^{\infty}(-1)^{r+1}\frac{x^{2r}}{(2r)!}
\end{align*}

Also,
\begin{align*}
x\sin x
&=x\sum_{r=0}^{\infty}(-1)^r\frac{x^{2r+1}}{(2r+1)!}\\
&=\sum_{r=0}^{\infty}(-1)^r\frac{x^{2r+2}}{(2r+1)!}
\end{align*}

Re-index this sum by writing \(r-1\) instead of \(r\), so that it is also in powers of \(x^{2r}\):
\[
x\sin x=\sum_{r=1}^{\infty}(-1)^{r-1}\frac{x^{2r}}{(2r-1)!}
\]

Hence
\[
2(1-\cos x)-x\sin x
=
\sum_{r=1}^{\infty}
\left(
\frac{2(-1)^{r+1}}{(2r)!}
-
\frac{(-1)^{r-1}}{(2r-1)!}
\right)x^{2r}
\]

So the coefficient of \(x^{2r}\) is
\begin{align*}
\frac{2(-1)^{r+1}}{(2r)!}-\frac{(-1)^{r-1}}{(2r-1)!}
&=(-1)^r\left(-\frac{2}{(2r)!}+\frac{1}{(2r-1)!}\right)\\
&=(-1)^r\left(-\frac{2}{(2r)!}+\frac{2r}{(2r)!}\right)\\
&=(-1)^r\frac{2(r-1)}{(2r)!}
\end{align*}

Therefore the general term of the Maclaurin series is
\[
(-1)^r\frac{2(r-1)}{(2r)!}x^{2r}
\]

as required.

  \item Using part (a), the first non-zero term in the numerator comes from \(r=2\):
\begin{align*}
2(1-\cos x)-x\sin x
&=\frac{2(2-1)}{4!}x^4+\cdots\\
&=\frac{x^4}{12}+\cdots
\end{align*}

Now expand \(1-\cos x\):
\begin{align*}
1-\cos x
&=\frac{x^2}{2!}-\frac{x^4}{4!}+\cdots\\
&=\frac{x^2}{2}-\frac{x^4}{24}+\cdots
\end{align*}

Squaring gives
\begin{align*}
(1-\cos x)^2
&=\left(\frac{x^2}{2}-\frac{x^4}{24}+\cdots\right)^2\\
&=\frac{x^4}{4}+\cdots
\end{align*}

Therefore
\[
\frac{2(1-\cos x)-x\sin x}{(1-\cos x)^2}
=
\frac{\frac{x^4}{12}+\cdots}{\frac{x^4}{4}+\cdots}
\]
and so, comparing the leading terms as \(x\to 0\),
\[
\lim_{x\to 0}\left[\frac{2(1-\cos x)-x\sin x}{(1-\cos x)^2}\right]
=
\frac{1/12}{1/4}
=
\frac13
\]

So the limit is \(\frac13\).
\end{enumerate}
\end{tcolorbox}


\newpage

% ---- Question 4 ----
\questionitem
% [Source: _QBT___Solns__Maclaurin_Series, Q4]

Let $f(x)=\ln(1+\sin x)$. \\
\begin{enumerate}[label=\textbf{(\alph*)}]
  \item
  \begin{enumerate}[label=(\roman*)]
    \item Determine $f''(x)$. \marks{2}
    \item Determine the first two non-zero terms of the Maclaurin expansion for $f(x)$. \marks{3}
    \item By considering the first two non-zero terms of the Maclaurin expansion for $f(x)$, find an approximation to
    \[
    \int_0^{\frac{2}{5}} f(x)\,\mathrm{d}x
    \]
    Give your answer correct to $6$ decimal places. \marks{2} \\
  \end{enumerate}
\end{enumerate}

\vspace*{10pt}

\begin{tcolorbox}[colback=white, colframe=scarlet, title={\textbf{Solution}}, breakable, grow to left by=6mm, grow to right by=10mm]
\begin{enumerate}[label=\textbf{(\alph*)}]
  \item
  \begin{enumerate}[label=(\roman*)]
    \item
    Given
    \[
    f(x)=\ln(1+\sin x)
    \]
    first differentiate using
    \[
    \frac{\mathrm{d}}{\mathrm{d}x}\bigl(\ln u\bigr)=\frac{1}{u}\frac{\mathrm{d}u}{\mathrm{d}x}
    \]
    so
    \[
    f'(x)=\frac{\cos x}{1+\sin x}
    \]
    Differentiate again using the quotient rule:
    \begin{align*}
    f''(x)
    &=\frac{(-\sin x)(1+\sin x)-\cos x(\cos x)}{(1+\sin x)^2} \\
    &=\frac{-\sin x-\sin^2x-\cos^2x}{(1+\sin x)^2} \\
    &=\frac{-\sin x-1}{(1+\sin x)^2} \\
    &=-\frac{1}{1+\sin x}
    \end{align*}
    Hence
    \[
    f''(x)=-\frac{1}{1+\sin x}
    \]

    \item
    For the Maclaurin expansion, find the values at \(x=0\):
    \[
    f(0)=\ln(1+\sin 0)=\ln 1=0
    \]
    \[
    f'(0)=\frac{\cos 0}{1+\sin 0}=\frac{1}{1}=1
    \]
    \[
    f''(0)=-\frac{1}{1+\sin 0}=-1
    \]
    Using
    \[
    f(x)=f(0)+f'(0)x+\frac{f''(0)}{2!}x^2+\cdots
    \]
    gives
    \begin{align*}
    f(x)
    &=0+1\cdot x+\frac{-1}{2}x^2+\cdots \\
    &=x-\frac{x^2}{2}+\cdots
    \end{align*}
    So the first two non-zero terms are
    \[
    x-\frac{x^2}{2}
    \]

    \item
    Using the first two non-zero terms,
    \[
    f(x)\approx x-\frac{x^2}{2}
    \]
    Therefore
    \begin{align*}
    \int_0^{2/5} f(x)\,\mathrm{d}x
    &\approx \int_0^{2/5}\left(x-\frac{x^2}{2}\right)\,\mathrm{d}x \\
    &=\left[\frac{x^2}{2}-\frac{x^3}{6}\right]_0^{2/5} \\
    &=\frac{(2/5)^2}{2}-\frac{(2/5)^3}{6} \\
    &=\frac{2}{25}-\frac{4}{375} \\
    &=\frac{26}{375}
    \end{align*}
    Hence the required approximation is
    \[
    \frac{26}{375}\approx 0.069333
    \]
    correct to 6 decimal places.
  \end{enumerate}
\end{enumerate}
\end{tcolorbox}


\newpage

% ---- Question 5 ----
\questionitem
% [Source: _QBT___Solns__Maclaurin_Series, Q5]

\begin{enumerate}[label=\textbf{(\alph*)}]
  \item Given that \[y=\cosh^{-1}\left(\frac{3+x}{2+x}\right)\] show that \[(x+2)\sqrt{5+2x}\,\frac{\mathrm{d}y}{\mathrm{d}x}+1=0\] \marks{4} \\
  \item Hence find the first three terms in the Maclaurin series for $\cosh^{-1}\left(\frac{3+x}{2+x}\right)$ in the form
  \[
  a\ln\left(\frac{3+\sqrt{5}}{2}\right)+bx+cx^2
  \]
  where $a$, $b$ and $c$ are constants to be determined. \marks{5} \\
\end{enumerate}

\vspace*{10pt}

\begin{tcolorbox}[colback=white, colframe=scarlet, title={\textbf{Solution}}, breakable, grow to left by=6mm, grow to right by=10mm]
\begin{enumerate}[label=\textbf{(\alph*)}]
  \item Let
\[
y=\cosh^{-1}\left(\frac{x+3}{x+2}\right)
\]
Then
\[
\cosh y=\frac{x+3}{x+2}
\]

Differentiate implicitly with respect to \(x\):
\begin{align*}
\sinh y\;\frac{\mathrm{d}y}{\mathrm{d}x}
&=\frac{(x+2)(1)-(x+3)(1)}{(x+2)^2} \\
&=\frac{x+2-x-3}{(x+2)^2} \\
&=-\frac{1}{(x+2)^2}
\end{align*}

Now use
\[
\sinh^2 y=\cosh^2 y-1
\]
So
\begin{align*}
\sinh^2 y
&=\left(\frac{x+3}{x+2}\right)^2-1 \\
&=\frac{(x+3)^2-(x+2)^2}{(x+2)^2} \\
&=\frac{x^2+6x+9-(x^2+4x+4)}{(x+2)^2} \\
&=\frac{2x+5}{(x+2)^2}
\end{align*}

Since \(y=\cosh^{-1}(\cdots)\), we have \(y\ge 0\), so \(\sinh y\ge 0\). Also near \(x=0\), \(x+2>0\). Hence
\[
\sinh y=\frac{\sqrt{2x+5}}{x+2}
\]

Substituting into the differentiated equation:
\begin{align*}
\frac{\sqrt{2x+5}}{x+2}\frac{\mathrm{d}y}{\mathrm{d}x}
&=-\frac{1}{(x+2)^2}
\end{align*}
so
\[
\frac{\mathrm{d}y}{\mathrm{d}x}=-\frac{1}{(x+2)\sqrt{2x+5}}
\]

Therefore
\[
(x+2)\sqrt{5+2x}\,\frac{\mathrm{d}y}{\mathrm{d}x}+1=0
\]

So the given result is shown.

\item From part (a),
\[
(x+2)\sqrt{5+2x}\,\frac{\mathrm{d}y}{\mathrm{d}x}+1=0
\]

First divide through by \(\sqrt{5+2x}\):
\[
(x+2)\frac{\mathrm{d}y}{\mathrm{d}x}+\frac{1}{\sqrt{5+2x}}=0
\]

Differentiate implicitly with respect to \(x\):
\begin{align*}
\frac{\mathrm{d}}{\mathrm{d}x}\left((x+2)\frac{\mathrm{d}y}{\mathrm{d}x}\right)
+\frac{\mathrm{d}}{\mathrm{d}x}\left((5+2x)^{-1/2}\right)&=0 \\
(x+2)\frac{\mathrm{d}^2y}{\mathrm{d}x^2}
+\frac{\mathrm{d}y}{\mathrm{d}x}
-(5+2x)^{-3/2}&=0
\end{align*}

Now evaluate the required quantities at \(x=0\).

First,
\[
y(0)=\cosh^{-1}\left(\frac{3}{2}\right)
\]
Using
\[
\cosh^{-1}t=\ln\left(t+\sqrt{t^2-1}\right),
\]
we get
\begin{align*}
y(0)
&=\ln\left(\frac32+\sqrt{\frac94-1}\right) \\
&=\ln\left(\frac32+\sqrt{\frac54}\right) \\
&=\ln\left(\frac32+\frac{\sqrt5}{2}\right) \\
&=\ln\left(\frac{3+\sqrt5}{2}\right)
\end{align*}

Also from part (a),
\[
\frac{\mathrm{d}y}{\mathrm{d}x}
=-\frac{1}{(x+2)\sqrt{5+2x}}
\]
so
\[
\left.\frac{\mathrm{d}y}{\mathrm{d}x}\right|_{x=0}
=-\frac{1}{2\sqrt5}
\]

Substitute \(x=0\) into
\[
(x+2)\frac{\mathrm{d}^2y}{\mathrm{d}x^2}
+\frac{\mathrm{d}y}{\mathrm{d}x}
-(5+2x)^{-3/2}=0:
\]
\[
2\left.\frac{\mathrm{d}^2y}{\mathrm{d}x^2}\right|_{x=0}
+\left.\frac{\mathrm{d}y}{\mathrm{d}x}\right|_{x=0}
-5^{-3/2}=0
\]
Hence
\begin{align*}
2\left.\frac{\mathrm{d}^2y}{\mathrm{d}x^2}\right|_{x=0}
-\frac{1}{2\sqrt5}
-\frac{1}{5\sqrt5}&=0 \\
2\left.\frac{\mathrm{d}^2y}{\mathrm{d}x^2}\right|_{x=0}
&=\frac{1}{2\sqrt5}+\frac{1}{5\sqrt5} \\
&=\frac{7}{10\sqrt5}
\end{align*}
so
\[
\left.\frac{\mathrm{d}^2y}{\mathrm{d}x^2}\right|_{x=0}
=\frac{7}{20\sqrt5}
\]

Therefore the Maclaurin series is
\[
y
=
y(0)
+\left.\frac{\mathrm{d}y}{\mathrm{d}x}\right|_{x=0}x
+\frac{1}{2}\left.\frac{\mathrm{d}^2y}{\mathrm{d}x^2}\right|_{x=0}x^2
+\cdots
\]
Thus
\[
\cosh^{-1}\left(\frac{3+x}{2+x}\right)
=
\ln\left(\frac{3+\sqrt5}{2}\right)
-\frac{x}{2\sqrt5}
+\frac{7x^2}{40\sqrt5}
+\cdots
\]

So
\[
a=1,\qquad b=-\frac{1}{2\sqrt5},\qquad c=\frac{7}{40\sqrt5}
\]
\end{enumerate}
\end{tcolorbox}


\newpage

% ---- Question 6 ----
\questionitem
% [Source: _QBT___Solns__Maclaurin_Series, Q6]

\[ y=(1+\cosh x)^n \qquad n\geq 2\] \\
\begin{enumerate}[label=\textbf{(\alph*)}]
  \item
  \begin{enumerate}[label=(\roman*)]
    \item Show that
    \[
    \frac{\mathrm{d}^2 y}{\mathrm{d}x^2}=n^2(1+\cosh x)^n-n(2n-1)(1+\cosh x)^{n-1}
    \]
    \marks[-20pt]{4}
    \item Determine an expression for
    \[
    \frac{\mathrm{d}^4 y}{\mathrm{d}x^4}
    \]
    \marks[-20pt]{2}
  \end{enumerate}
  \item Hence determine the first three non-zero terms of the Maclaurin series for $y$, giving each coefficient in simplest form. \marks{2} \\
\end{enumerate}

\vspace*{10pt}

\begin{tcolorbox}[colback=white, colframe=scarlet, title={\textbf{Solution}}, breakable, grow to left by=6mm, grow to right by=10mm]
\begin{enumerate}[label=\textbf{(\alph*)}]
  \item
  \begin{enumerate}[label=\textbf{(\roman*)}]
    \item Let
    \[
    y=(1+\cosh x)^n
    \]

    Differentiate once:
    \[
    \frac{\mathrm{d}y}{\mathrm{d}x}=n(1+\cosh x)^{n-1}\sinh x
    \]

    Differentiate again using the product rule:
    \begin{align*}
    \frac{\mathrm{d}^2y}{\mathrm{d}x^2}
    &=n\left((n-1)(1+\cosh x)^{n-2}\sinh^2 x+(1+\cosh x)^{n-1}\cosh x\right) \\
    &=n(n-1)(1+\cosh x)^{n-2}\sinh^2 x+n(1+\cosh x)^{n-1}\cosh x
    \end{align*}

    Now use
    \[
    \sinh^2 x=\cosh^2 x-1=(\cosh x-1)(\cosh x+1)
    \]

    Since $\cosh x+1=1+\cosh x$,
    \begin{align*}
    \frac{\mathrm{d}^2y}{\mathrm{d}x^2}
    &=n(n-1)(1+\cosh x)^{n-2}(\cosh x-1)(1+\cosh x)+n(1+\cosh x)^{n-1}\cosh x \\
    &=n(n-1)(1+\cosh x)^{n-1}(\cosh x-1)+n(1+\cosh x)^{n-1}\cosh x \\
    &=n(1+\cosh x)^{n-1}\left((n-1)(\cosh x-1)+\cosh x\right) \\
    &=n(1+\cosh x)^{n-1}\left(n\cosh x-(n-1)\right)
    \end{align*}

    Write
    \[
    n\cosh x=n(1+\cosh x)-n
    \]

    Then
    \[
    n\cosh x-(n-1)=n(1+\cosh x)-(2n-1)
    \]

    So
    \begin{align*}
    \frac{\mathrm{d}^2y}{\mathrm{d}x^2}
    &=n(1+\cosh x)^{n-1}\left(n(1+\cosh x)-(2n-1)\right) \\
    &=n^2(1+\cosh x)^n-n(2n-1)(1+\cosh x)^{n-1}
    \end{align*}

    Hence
    \[
    \frac{\mathrm{d}^2 y}{\mathrm{d}x^2} = n^2(1+\cosh x)^n - n(2n-1)(1+\cosh x)^{n-1}
    \]

    \item From part (i),
    \[
    y''=n^2(1+\cosh x)^n-n(2n-1)(1+\cosh x)^{n-1}
    \]

    Differentiate twice again. Using part (i) with powers $n$ and $n-1$:

    \[
    \frac{\mathrm{d}^2}{\mathrm{d}x^2}(1+\cosh x)^n
    =n^2(1+\cosh x)^n-n(2n-1)(1+\cosh x)^{n-1}
    \]

    and
    \[
    \frac{\mathrm{d}^2}{\mathrm{d}x^2}(1+\cosh x)^{n-1}
    =(n-1)^2(1+\cosh x)^{n-1}-(n-1)(2n-3)(1+\cosh x)^{n-2}
    \]

    Therefore
    \begin{align*}
    \frac{\mathrm{d}^4y}{\mathrm{d}x^4}
    &=n^2\frac{\mathrm{d}^2}{\mathrm{d}x^2}(1+\cosh x)^n
    -n(2n-1)\frac{\mathrm{d}^2}{\mathrm{d}x^2}(1+\cosh x)^{n-1} \\
    &=n^2\left[n^2(1+\cosh x)^n-n(2n-1)(1+\cosh x)^{n-1}\right] \\
    &\quad -n(2n-1)\left[(n-1)^2(1+\cosh x)^{n-1}-(n-1)(2n-3)(1+\cosh x)^{n-2}\right] \\
    &=n^4(1+\cosh x)^n-n(2n-1)\left(n^2+(n-1)^2\right)(1+\cosh x)^{n-1} \\
    &\quad +n(n-1)(2n-1)(2n-3)(1+\cosh x)^{n-2}
    \end{align*}

    Since
    \[
    n^2+(n-1)^2=2n^2-2n+1
    \]

    we get
    \[
    \frac{\mathrm{d}^4 y}{\mathrm{d}x^4}
    =n^4(1+\cosh x)^n-n(2n-1)(2n^2-2n+1)(1+\cosh x)^{n-1}
    +n(n-1)(2n-1)(2n-3)(1+\cosh x)^{n-2}
    \]
  \end{enumerate}

  \item Since $\cosh x$ is an even function, $y=(1+\cosh x)^n$ is also even. So its Maclaurin series contains only even powers, and
  \[
  y'(0)=0,\qquad y'''(0)=0
  \]

  Also,
  \[
  y(0)=(1+\cosh 0)^n=(1+1)^n=2^n
  \]

  From part (a)(i),
  \begin{align*}
  y''(0)
  &=n^2(2^n)-n(2n-1)(2^{\,n-1}) \\
  &=n2^{\,n-1}\left(2n-(2n-1)\right) \\
  &=n2^{\,n-1}
  \end{align*}

  From part (a)(ii),
  \begin{align*}
  y''''(0)
  &=n^4 2^n-n(2n-1)(2n^2-2n+1)2^{\,n-1}+n(n-1)(2n-1)(2n-3)2^{\,n-2} \\
  &=n2^{\,n-2}\left[4n^3-2(2n-1)(2n^2-2n+1)+(n-1)(2n-1)(2n-3)\right]
  \end{align*}

  Expand the bracket:
  \begin{align*}
  2(2n-1)(2n^2-2n+1)&=8n^3-12n^2+8n-2 \\
  (n-1)(2n-1)(2n-3)&=4n^3-12n^2+11n-3
  \end{align*}

  Hence
  \begin{align*}
  y''''(0)
  &=n2^{\,n-2}\left[4n^3-(8n^3-12n^2+8n-2)+(4n^3-12n^2+11n-3)\right] \\
  &=n2^{\,n-2}(3n-1)
  \end{align*}

  Using
  \[
  y=y(0)+\frac{y''(0)}{2!}x^2+\frac{y''''(0)}{4!}x^4+\cdots
  \]

  gives
  \begin{align*}
  y
  &=2^n+\frac{n2^{\,n-1}}{2}x^2+\frac{n(3n-1)2^{\,n-2}}{24}x^4+\cdots \\
  &=2^n+n2^{\,n-2}x^2+\frac{n(3n-1)2^{\,n-2}}{24}x^4+\cdots
  \end{align*}

  So the first three non-zero terms are
  \[
  2^n+n2^{\,n-2}x^2+\frac{n(3n-1)2^{\,n-2}}{24}x^4
  \]
\end{enumerate}
\end{tcolorbox}


\newpage

% ---- Question 7 ----
\questionitem
% [Source: _QBT___Solns__Maclaurin_Series, Q7]

\begin{enumerate}[label=(\roman*)]
  \item Prove that, if $y=\arctan x$, then
  \[
  \frac{\mathrm{d}y}{\mathrm{d}x}=\frac{1}{1+x^2}
  \]
  \marks[-20pt]{3}
  \item Find the Maclaurin series for $\arctan x$, up to and including the term in $x^5$. \marks{4} \\
  \item Use the result of part (ii) and the Maclaurin series for $\ln(1+t)$ to find the Maclaurin series for $\ln(1+\arctan x)$, up to and including the term in $x^4$. \marks{5} \\
\end{enumerate}

\vspace*{10pt}

\begin{tcolorbox}[colback=white, colframe=scarlet, title={\textbf{Solution}}, breakable, grow to left by=6mm, grow to right by=10mm]
\begin{enumerate}[label=\textbf{(\roman*)}]
  \item Let
  \[
  y=\arctan x
  \]
  so equivalently
  \[
  x=\tan y
  \]
  Differentiate implicitly with respect to $x$:
  \[
  \frac{\mathrm{d}}{\mathrm{d}x}(x)=\frac{\mathrm{d}}{\mathrm{d}x}(\tan y)
  \]
  \[
  1=\sec^2 y\frac{\mathrm{d}y}{\mathrm{d}x}
  \]
  Hence
  \[
  \frac{\mathrm{d}y}{\mathrm{d}x}=\frac{1}{\sec^2 y}
  \]
  Using the identity $\sec^2 y=1+\tan^2 y$, and $\tan y=x$, we get
  \[
  \frac{\mathrm{d}y}{\mathrm{d}x}=\frac{1}{1+\tan^2 y}=\frac{1}{1+x^2}
  \]
  Therefore
  \[
  \frac{\mathrm{d}}{\mathrm{d}x}(\arctan x)=\frac{1}{1+x^2}
  \]

  \item From part (i),
  \[
  \frac{\mathrm{d}}{\mathrm{d}x}(\arctan x)=\frac{1}{1+x^2}
  \]
  Using the geometric series
  \[
  \frac{1}{1-r}=1+r+r^2+r^3+\cdots \qquad (|r|<1)
  \]
  with $r=-x^2$, we have
  \[
  \frac{1}{1+x^2}=1-x^2+x^4-x^6+\cdots \qquad (|x|<1)
  \]
  So
  \[
  \frac{\mathrm{d}}{\mathrm{d}x}(\arctan x)=1-x^2+x^4-\cdots
  \]
  Integrating term by term,
  \[
  \arctan x=\int (1-x^2+x^4-\cdots)\,\mathrm{d}x
  \]
  \[
  \arctan x=x-\frac{x^3}{3}+\frac{x^5}{5}+\cdots +C
  \]
  To find $C$, put $x=0$:
  \[
  \tan^{-1}0=0
  \]
  so $C=0$.

  Therefore the Maclaurin series is
  \[
  \arctan x=x-\frac{x^3}{3}+\frac{x^5}{5}+\cdots
  \]

  \item Let
  \[
  t=\arctan x
  \]
  Then from part (ii),
  \[
  t=x-\frac{x^3}{3}+\cdots
  \]
  We use the Maclaurin series
  \[
  \ln(1+t)=t-\frac{t^2}{2}+\frac{t^3}{3}-\frac{t^4}{4}+\cdots
  \]

  Now find the powers of $t$ needed up to the term in $x^4$.

  First,
  \[
  t^2=\left(x-\frac{x^3}{3}+\cdots\right)^2
  \]
  \[
  t^2=x^2-\frac{2x^4}{3}+\cdots
  \]

  Next,
  \[
  t^3=\left(x-\frac{x^3}{3}+\cdots\right)^3=x^3+\cdots
  \]

  Also,
  \[
  t^4=\left(x-\frac{x^3}{3}+\cdots\right)^4=x^4+\cdots
  \]

  Substitute into the logarithm series:
  \begin{align*}
  \ln(1+\arctan x)
  &=\ln(1+t)\\
  &=t-\frac{t^2}{2}+\frac{t^3}{3}-\frac{t^4}{4}+\cdots\\
  &=\left(x-\frac{x^3}{3}+\cdots\right)
   -\frac12\left(x^2-\frac{2x^4}{3}+\cdots\right)
   +\frac13\left(x^3+\cdots\right)
   -\frac14\left(x^4+\cdots\right)
   +\cdots
  \end{align*}

  Now simplify:
  \begin{align*}
  \ln(1+\arctan x)
  &=x-\frac{x^3}{3}-\frac{x^2}{2}+\frac{x^4}{3}+\frac{x^3}{3}-\frac{x^4}{4}+\cdots\\
  &=x-\frac{x^2}{2}+\left(\frac13-\frac14\right)x^4+\cdots\\
  &=x-\frac{x^2}{2}+\frac{x^4}{12}+\cdots
  \end{align*}

  Therefore
  \[
  \ln(1+\arctan x)=x-\frac{x^2}{2}+\frac{x^4}{12}+\cdots
  \]
\end{enumerate}
\end{tcolorbox}


\newpage

% ---- Question 8 ----
\questionitem
% [Source: _QBT___Solns__Maclaurin_Series, Q8]

\begin{enumerate}[label=\textbf{(\alph*)}]
  \item The function $f$ is defined as $f(x)=\operatorname{sech}^2 x$. \\
  \begin{enumerate}[label=(\roman*)]
    \item Show that $f^{(4)}(0)=16$. \marks{4}
    \item Hence find the first three non-zero terms of the Maclaurin series for $f(x)=\operatorname{sech}^2 x$. \marks{2} \\
  \end{enumerate}
  \item Prove that
  \[
  \lim_{x\to 0}\left(\frac{\operatorname{sech}^2 x-\cos(\sqrt{2}\,x)}{x^4}\right)=\frac{1}{2}
  \]
  \marks[-30pt]{4}
\end{enumerate}

\vspace*{10pt}

\begin{tcolorbox}[colback=white, colframe=scarlet, title={\textbf{Solution}}, breakable, grow to left by=6mm, grow to right by=10mm]
\begin{enumerate}[label=\textbf{(\alph*)}]
  \item
  \begin{enumerate}[label=(\roman*)]
    \item Let
    \[
    f(x)=\operatorname{sech}^2 x
    \]
    Using
    \[
    \frac{\mathrm{d}}{\mathrm{d}x}(\operatorname{sech}x)=-\operatorname{sech}x\tanh x
    \]
    we get
    \[
    f'(x)=2\operatorname{sech}x\bigl(-\operatorname{sech}x\tanh x\bigr)=-2\operatorname{sech}^2x\tanh x=-2f\tanh x
    \]

    Differentiate again:
    \begin{align*}
    f''(x)
    &= -2f'\tanh x-2f\operatorname{sech}^2x \\
    &= -2f'\tanh x-2f^2
    \end{align*}

    Since
    \[
    \tanh^2x=1-\operatorname{sech}^2x=1-f
    \]
    and \(f'=-2f\tanh x\),
    \begin{align*}
    f''(x)
    &= -2(-2f\tanh x)\tanh x-2f^2 \\
    &= 4f\tanh^2x-2f^2 \\
    &= 4f(1-f)-2f^2 \\
    &= 4f-6f^2
    \end{align*}

    Now evaluate at \(x=0\). Since \(\operatorname{sech}0=1\) and \(\tanh 0=0\),
    \[
    f(0)=1
    \]
    \[
    f'(0)=-2f(0)\tanh 0=0
    \]
    \[
    f''(0)=4f(0)-6f(0)^2=4-6=-2
    \]

    Differentiate \(f''=4f-6f^2\):
    \[
    f'''=4f'-12ff'
    \]
    so
    \[
    f'''(0)=4f'(0)-12f(0)f'(0)=0
    \]

    Differentiate once more:
    \[
    f^{(4)}=4f''-12\bigl((f')^2+ff''\bigr)
    \]
    Hence
    \begin{align*}
    f^{(4)}(0)
    &= 4f''(0)-12\bigl((f'(0))^2+f(0)f''(0)\bigr) \\
    &= 4(-2)-12\bigl(0+1(-2)\bigr) \\
    &= -8+24 \\
    &= 16
    \end{align*}

    Therefore
    \[
    f^{(4)}(0)=16
    \]

    \item The Maclaurin series for \(f(x)\) is
    \[
    f(x)=f(0)+f'(0)x+\frac{f''(0)}{2!}x^2+\frac{f'''(0)}{3!}x^3+\frac{f^{(4)}(0)}{4!}x^4+\cdots
    \]

    Substituting the values found in part (i),
    \begin{align*}
    f(x)
    &= 1+0x+\frac{-2}{2!}x^2+\frac{0}{3!}x^3+\frac{16}{4!}x^4+\cdots \\
    &= 1-x^2+\frac{16}{24}x^4+\cdots \\
    &= 1-x^2+\frac{2}{3}x^4+\cdots
    \end{align*}

    So the first three non-zero terms are
    \[
    \operatorname{sech}^2x=1-x^2+\frac{2}{3}x^4+\cdots
    \]
  \end{enumerate}

  \item From part (a),
  \[
  \operatorname{sech}^2x=1-x^2+\frac{2}{3}x^4+\cdots
  \]

  Also, from the standard Maclaurin series,
  \[
  \cos t=1-\frac{t^2}{2!}+\frac{t^4}{4!}+\cdots
  \]
  Put \(t=\sqrt{2}\,x\):
  \begin{align*}
  \cos(\sqrt{2}\,x)
  &= 1-\frac{(\sqrt{2}\,x)^2}{2!}+\frac{(\sqrt{2}\,x)^4}{4!}+\cdots \\
  &= 1-\frac{2x^2}{2}+\frac{4x^4}{24}+\cdots \\
  &= 1-x^2+\frac{1}{6}x^4+\cdots
  \end{align*}

  Therefore
  \begin{align*}
  \operatorname{sech}^2x-\cos(\sqrt{2}\,x)
  &= \left(1-x^2+\frac{2}{3}x^4+\cdots\right)-\left(1-x^2+\frac{1}{6}x^4+\cdots\right) \\
  &= \left(\frac{2}{3}-\frac{1}{6}\right)x^4+\cdots \\
  &= \frac{1}{2}x^4+\cdots
  \end{align*}

  So
  \[
  \frac{\operatorname{sech}^2x-\cos(\sqrt{2}\,x)}{x^4}
  = \frac{1}{2}+\text{terms involving positive powers of }x
  \]

  As \(x\to 0\), those remaining terms tend to \(0\). Hence
  \[
  \lim_{x\to 0}\left(\frac{\operatorname{sech}^2 x-\cos(\sqrt{2}\,x)}{x^4}\right)=\frac{1}{2}
  \]
\end{enumerate}
\end{tcolorbox}


\newpage

% ---- Question 9 ----
\questionitem
% [Source: _QBT___Solns__Maclaurin_Series, Q9]

Using the Maclaurin series for $\sin x$, show that, for small values of $x$,
\[
\frac{x}{\sin x} \approx 1 + ax^2 + bx^4 + cx^6
\]
where the values of $a$, $b$ and $c$ are to be given in exact form. \marks{5}

\vspace*{10pt}

\begin{tcolorbox}[colback=white, colframe=scarlet, title={\textbf{Solution}}, breakable, grow to left by=6mm, grow to right by=10mm]
Using the Maclaurin series,
\[
\sin x=x-\frac{x^3}{3!}+\frac{x^5}{5!}-\frac{x^7}{7!}+\cdots
\]
so
\[
\sin x=x-\frac{x^3}{6}+\frac{x^5}{120}-\frac{x^7}{5040}+\cdots
\]

Factor out $x$:
\[
\sin x=x\left(1-\frac{x^2}{6}+\frac{x^4}{120}-\frac{x^6}{5040}+\cdots\right)
\]

Hence
\[
\frac{x}{\sin x}
=\frac{1}{1-\frac{x^2}{6}+\frac{x^4}{120}-\frac{x^6}{5040}+\cdots}
\]

We are told that, for small $x$,
\[
\frac{x}{\sin x}\approx 1+ax^2+bx^4+cx^6
\]

So this expression must satisfy
\[
(1+ax^2+bx^4+cx^6)\left(1-\frac{x^2}{6}+\frac{x^4}{120}-\frac{x^6}{5040}\right)=1+\cdots
\]

Now expand and collect powers of $x$:
\begin{align*}
&(1+ax^2+bx^4+cx^6)\left(1-\frac{x^2}{6}+\frac{x^4}{120}-\frac{x^6}{5040}\right) \\
&=1+\left(a-\frac16\right)x^2+\left(b-\frac{a}{6}+\frac1{120}\right)x^4
+\left(c-\frac{b}{6}+\frac{a}{120}-\frac1{5040}\right)x^6+\cdots
\end{align*}

For this to equal \(1\) up to and including the term in \(x^6\), the coefficients of \(x^2\), \(x^4\) and \(x^6\) must each be \(0\).

So:
\[
a-\frac16=0
\qquad\Rightarrow\qquad
a=\frac16
\]

\[
b-\frac{a}{6}+\frac1{120}=0
\qquad\Rightarrow\qquad
b=\frac{a}{6}-\frac1{120}
\]
Substituting $a=\frac16$,
\[
b=\frac{1}{36}-\frac{1}{120}
=\frac{7}{360}
\]

\[
c-\frac{b}{6}+\frac{a}{120}-\frac1{5040}=0
\qquad\Rightarrow\qquad
c=\frac{b}{6}-\frac{a}{120}+\frac1{5040}
\]
Substituting $a=\frac16$ and $b=\frac7{360}$,
\begin{align*}
c&=\frac{7}{2160}-\frac{1}{720}+\frac{1}{5040}\\
&=\frac{49-21+3}{15120}\\
&=\frac{31}{15120}
\end{align*}

Therefore,
\[
a=\frac16,\qquad b=\frac7{360},\qquad c=\frac{31}{15120}
\]

and
\[
\frac{x}{\sin x}\approx 1+\frac{x^2}{6}+\frac{7x^4}{360}+\frac{31x^6}{15120}
\]
for small values of $x$.
\end{tcolorbox}


\newpage

% ---- Question 10 ----
\questionitem
% [Source: _QBT___Solns__Maclaurin_Series, Q10]

\begin{enumerate}[label=\textbf{(\alph*)}]
  \item Use the Maclaurin series expansion for $\ln(1+x)$ to show that the first three non-zero terms of the Maclaurin series expansion of
  \[
  \ln\left(\frac{1+2x}{1-x}\right)
  \]
  are
  \[
  3x-\frac{3}{2}x^2+3x^3
  \]
  \marks[-20pt]{2}
  \item Sam attempts to use the series expansion found in part (a) to find an approximation for $\ln 6$. \\

  Sam's incorrect working is shown below.
\begin{tcolorbox}[enhanced,boxrule=0.4pt,colback=white,arc=0mm,outer arc=0mm]
  \[
  \begin{aligned}
  \text{Let } \frac{1+2x}{1-x} &= 6\\
  1+2x &= 6-6x\\
  8x &= 5\\
  x &= \frac{5}{8} 
  \end{aligned}
  \] \\
  \[
  \begin{aligned}
  \text{So } \ln 6 &\approx 3\left(\frac{5}{8}\right)-\frac{3}{2}\left(\frac{5}{8}\right)^2+3\left(\frac{5}{8}\right)^3\\
  &\approx 2.02
  \end{aligned}
  \]
\end{tcolorbox}
  Explain the error in Sam's working. \marks{2} \\
  \item Use $x=\frac{1}{4}$ in the series expansion found in part (a) to find an approximation for $\ln 8$. \\

  Fully justify your answer. \marks{3} \\
\end{enumerate}

\vspace*{10pt}

\begin{tcolorbox}[colback=white, colframe=scarlet, title={\textbf{Solution}}, breakable, grow to left by=6mm, grow to right by=10mm]
\begin{enumerate}[label=\textbf{(\alph*)}]
  \item Using
\[
\ln\left(\frac{1+2x}{1-x}\right)=\ln(1+2x)-\ln(1-x)
\]

and the Maclaurin series
\[
\ln(1+u)=u-\frac{u^2}{2}+\frac{u^3}{3}+\cdots
\]

we have, with $u=2x$,
\[
\ln(1+2x)=2x-\frac{(2x)^2}{2}+\frac{(2x)^3}{3}+\cdots
=2x-2x^2+\frac{8}{3}x^3+\cdots
\]

Also, with $u=-x$,
\[
\ln(1-x)=(-x)-\frac{(-x)^2}{2}+\frac{(-x)^3}{3}+\cdots
=-x-\frac12x^2-\frac13x^3+\cdots
\]

So
\begin{align*}
\ln\left(\frac{1+2x}{1-x}\right)
&=\left(2x-2x^2+\frac{8}{3}x^3+\cdots\right)-\left(-x-\frac12x^2-\frac13x^3+\cdots\right) \\
&=3x-\frac32x^2+3x^3+\cdots
\end{align*}

Hence the first three non-zero terms are
\[
3x-\frac32x^2+3x^3
\]

  \item Sam’s algebra for finding $x$ is correct, since
\[
\frac{1+2x}{1-x}=6 \quad \Rightarrow \quad x=\frac58
\]

The error is that the series from part (a) is not valid for this value of $x$.

For $\ln(1+2x)$, the Maclaurin series for $\ln(1+u)$ requires
\[
-1<u\leq 1
\]
so here
\[
-1<2x\leq 1 \quad \Rightarrow \quad -\frac12<x\leq \frac12
\]

Therefore the series for
\[
\ln\left(\frac{1+2x}{1-x}\right)
\]
is only valid for
\[
-\frac12<x\leq \frac12
\]

But
\[
\frac58>\frac12
\]

So Sam substituted a value of $x$ outside the interval of convergence, and the approximation is invalid.

  \item Put $x=\frac14$ into the expression inside the logarithm:
\[
\frac{1+2x}{1-x}
=\frac{1+2\left(\frac14\right)}{1-\frac14}
=\frac{1+\frac12}{\frac34}
=\frac{\frac32}{\frac34}
=2
\]

So the series from part (a) gives an approximation for $\ln 2$:
\begin{align*}
\ln 2
&\approx 3\left(\frac14\right)-\frac32\left(\frac14\right)^2+3\left(\frac14\right)^3 \\
&=\frac34-\frac{3}{32}+\frac{3}{64} \\
&=\frac{48-6+3}{64} \\
&=\frac{45}{64}
\end{align*}

Now
\[
\ln 8=\ln(2^3)=3\ln 2
\]

Hence
\[
\ln 8 \approx 3\times \frac{45}{64}=\frac{135}{64}=2.109375
\]

So
\[
\ln 8 \approx \frac{135}{64} \approx 2.11
\]

This is justified because $x=\frac14$ lies in the valid interval
\[
-\frac12<x\leq \frac12
\]

Also, the next term in the series is found by continuing part (a):
\[
\ln(1+2x)=2x-2x^2+\frac83x^3-4x^4+\cdots
\]
\[
\ln(1-x)=-x-\frac12x^2-\frac13x^3-\frac14x^4+\cdots
\]
so
\[
\ln\left(\frac{1+2x}{1-x}\right)=3x-\frac32x^2+3x^3-\frac{15}{4}x^4+\cdots
\]

At $x=\frac14$ this becomes
\[
\ln 2=\frac34-\frac{3}{32}+\frac{3}{64}-\frac{15}{1024}+\cdots
\]

This is an alternating series with decreasing terms, so the error is less than the first omitted term:
\[
\left|\text{error in } \ln 2\right|<\frac{15}{1024}
\]

Therefore
\[
\left|\text{error in } \ln 8\right|<3\cdot \frac{15}{1024}=\frac{45}{1024}\approx 0.0439
\]

Since the next omitted term is negative, $\frac{135}{64}$ is an overestimate.

Hence the approximation is
\[
\ln 8 \approx \frac{135}{64} \approx 2.11
\]
  \end{enumerate}
\end{tcolorbox}


\newpage

% ---- Question 11 ----
\questionitem
% [Source: _QBT___Solns__Maclaurin_Series, Q11]

\begin{enumerate}[label=\textbf{(\alph*)}]
  \item By using an appropriate Maclaurin series prove that if $0 < x < 1$ then $-\ln(1-x) > x$. \marks{2} \\
  \item Hence, by using a suitable substitution, deduce that $\ln t > 1 - \frac{1}{t}$ for $t > 1$. \marks{1} \\
  \item Using the inequality in part (b), and by making a suitable choice for $t$, determine which is greater, \[\left(\frac{5}{4}\right)^5 \quad \text{or}\quad e\] \marks{3} \\
\end{enumerate}

\vspace*{10pt}

\begin{tcolorbox}[colback=white, colframe=scarlet, title={\textbf{Solution}}, breakable, grow to left by=6mm, grow to right by=10mm]
\begin{enumerate}[label=\textbf{(\alph*)}]
  \item Using the Maclaurin series
  \[
  \ln(1+u)=u-\frac{u^2}{2}+\frac{u^3}{3}-\frac{u^4}{4}+\cdots \qquad (|u|<1)
  \]
  put $u=-x$. Then, for $0<x<1$,
  \begin{align*}
  \ln(1-x)&=-x-\frac{x^2}{2}-\frac{x^3}{3}-\frac{x^4}{4}-\cdots \\
  \therefore\quad -\ln(1-x)&=x+\frac{x^2}{2}+\frac{x^3}{3}+\frac{x^4}{4}+\cdots
  \end{align*}
  Since $0<x<1$, every term after the first is positive, so
  \[
  -\ln(1-x)>x
  \]
  Hence proved.

  \item Let
  \[
  x=1-\frac1t=\frac{t-1}{t}
  \]
  If $t>1$, then $0<x<1$, so part (a) applies:
  \[
  -\ln(1-x)>x
  \]
  Substituting $x=1-\frac1t$ gives
  \begin{align*}
  -\ln\!\left(1-\left(1-\frac1t\right)\right)&>1-\frac1t \\
  -\ln\!\left(\frac1t\right)&>1-\frac1t \\
  \ln t&>1-\frac1t
  \end{align*}
  So, for $t>1$,
  \[
  \ln t>1-\frac1t
  \]

  \item Choose $t=\dfrac54$ in part (b). Since $\dfrac54>1$,
  \[
  \ln\left(\frac54\right)>1-\frac{1}{5/4}=1-\frac45=\frac15
  \]
  Multiply both sides by $5$:
  \[
  5\ln\left(\frac54\right)>1
  \]
  But $1=\ln e$, so
  \[
  \ln\left(\left(\frac54\right)^5\right)>\ln e
  \]
  Since $\ln x$ is an increasing function,
  \[
  \left(\frac54\right)^5>e
  \]
  Therefore, $\left(\dfrac54\right)^5$ is greater than $e$.
\end{enumerate}
\end{tcolorbox}


\newpage

% ---- Question 12 ----
\questionitem
% [Source: _QBT___Solns__Maclaurin_Series, Q12]

\begin{enumerate}[label=\textbf{(\alph*)}]
  \item Write down the first three non-zero terms of the Maclaurin series for $\ln\left(\dfrac{1+2x}{1-x}\right)$. \marks{2} \\
  \item Use these three terms to show that
  \[
  \ln\left(\dfrac{10}{7}\right) \approx \dfrac{183}{512}
  \]
  \marks[-20pt]{2}
  \item Dana uses the same first three terms of the series to approximate $\ln 10$ and gets an answer of $2.67$, correct to $3$ significant figures. However, $\ln 10 = 2.30$ correct to $3$ significant figures. \\

  Explain Dana's error. \marks{2} \\
\end{enumerate}

\vspace*{10pt}

\begin{tcolorbox}[colback=white, colframe=scarlet, title={\textbf{Solution}}, breakable, grow to left by=6mm, grow to right by=10mm]
\begin{enumerate}[label=\textbf{(\alph*)}]
  \item Using
  \[
  \ln(1+u)=u-\frac{u^2}{2}+\frac{u^3}{3}+\cdots
  \]
  we have
  \[
  \ln\left(\frac{1+2x}{1-x}\right)=\ln(1+2x)-\ln(1-x)
  \]
  Now expand each logarithm:
  \[
  \ln(1+2x)=2x-\frac{(2x)^2}{2}+\frac{(2x)^3}{3}+\cdots
  =2x-2x^2+\frac{8x^3}{3}+\cdots
  \]
  \[
  \ln(1-x)=-x-\frac{x^2}{2}-\frac{x^3}{3}+\cdots
  \]
  So
  \begin{align*}
  \ln\left(\frac{1+2x}{1-x}\right)
  &=\left(2x-2x^2+\frac{8x^3}{3}\right)-\left(-x-\frac{x^2}{2}-\frac{x^3}{3}\right)+\cdots \\
  &=3x-\frac{3x^2}{2}+3x^3+\cdots
  \end{align*}
  The first three non-zero terms are
  \[
  3x-\frac{3x^2}{2}+3x^3
  \]

  \item To use the series for \(\ln\left(\dfrac{10}{7}\right)\), first find \(x\) such that
  \[
  \frac{1+2x}{1-x}=\frac{10}{7}
  \]
  Hence
  \[
  7(1+2x)=10(1-x)
  \]
  \[
  7+14x=10-10x
  \]
  \[
  24x=3 \implies x=\frac18
  \]
  Substitute \(x=\dfrac18\) into the first three terms:
  \begin{align*}
  \ln\left(\frac{10}{7}\right)
  &\approx 3\left(\frac18\right)-\frac32\left(\frac18\right)^2+3\left(\frac18\right)^3 \\
  &=\frac38-\frac{3}{128}+\frac{3}{512} \\
  &=\frac{192}{512}-\frac{12}{512}+\frac{3}{512} \\
  &=\frac{183}{512}
  \end{align*}
  So
  \[
  \ln\left(\frac{10}{7}\right)\approx \frac{183}{512}
  \]

  \item For \(\ln 10\), Dana must have used
  \[
  \frac{1+2x}{1-x}=10
  \]
  so
  \[
  1+2x=10-10x \implies 12x=9 \implies x=\frac34
  \]
  However, the series in part (a) comes from the Maclaurin expansions of \(\ln(1+2x)\) and \(\ln(1-x)\).

  The expansion for \(\ln(1+2x)\) is only valid when
  \[
  |2x|<1 \implies |x|<\frac12
  \]
  Since \(x=\dfrac34\), this is outside the interval where the series converges.

  So Dana's error was using the first three terms of the Maclaurin series at a value of \(x\) for which the series is not valid. That is why her answer \(2.67\) is not reliable.
\end{enumerate}
\end{tcolorbox}


\newpage

% ---- Question 13 ----
\questionitem
% [Source: _QBT___Solns__Maclaurin_Series, Q13]

\begin{enumerate}[label=\textbf{(\alph*)}]
  \item Use the Maclaurin series expansions for $\sin x$ and $\cos x$ to determine the series expansion of
  \[
  \left(\frac{\sin(x/2)}{x/2}\right)\cos\left(\frac{x}{3}\right)
  \]
  in ascending powers of $x$, up to and including the term in $x^4$. \\

  Give each term in simplest form. \marks{3} \\
  \item Use the answer to part (a) and calculus to find an approximation, to $5$ decimal places, for
  \[
  \int_{\pi/6}^{2\pi/3} \left( \frac{1}{x}\left(\frac{\sin(x/2)}{x/2}\right)\cos\left(\frac{x}{3}\right) \right) \, \mathrm{d}x
  \]
  \marks[-30pt]{3}
  \item Use the integration function on your calculator to evaluate
  \[
  \int_{\pi/6}^{2\pi/3} \left( \frac{1}{x}\left(\frac{\sin(x/2)}{x/2}\right)\cos\left(\frac{x}{3}\right) \right) \, \mathrm{d}x
  \]

  Give your answer to $5$ decimal places. \marks{1} \\
  \item Assuming that the calculator answer in part (c) is accurate to $5$ decimal places, comment on the accuracy of the approximation found in part (b). \marks{1} \\
\end{enumerate}

\vspace*{10pt}

\begin{tcolorbox}[colback=white, colframe=scarlet, title={\textbf{Solution}}, breakable, grow to left by=6mm, grow to right by=10mm]
\begin{enumerate}[label=\textbf{(\alph*)}]
  \item Using the standard Maclaurin series,
\[
\sin u=u-\frac{u^3}{3!}+\frac{u^5}{5!}+\cdots,
\qquad
\cos u=1-\frac{u^2}{2!}+\frac{u^4}{4!}+\cdots
\]

With \(u=\dfrac{x}{2}\),
\begin{align*}
\sin\left(\frac{x}{2}\right)
&=\frac{x}{2}-\frac{1}{3!}\left(\frac{x}{2}\right)^3+\frac{1}{5!}\left(\frac{x}{2}\right)^5+\cdots \\
&=\frac{x}{2}-\frac{x^3}{48}+\frac{x^5}{3840}+\cdots
\end{align*}
So
\[
\frac{\sin(x/2)}{x/2}
=1-\frac{x^2}{24}+\frac{x^4}{1920}+\cdots
\]

Also, with \(u=\dfrac{x}{3}\),
\begin{align*}
\cos\left(\frac{x}{3}\right)
&=1-\frac{1}{2!}\left(\frac{x}{3}\right)^2+\frac{1}{4!}\left(\frac{x}{3}\right)^4+\cdots \\
&=1-\frac{x^2}{18}+\frac{x^4}{1944}+\cdots
\end{align*}

Now multiply the two series, keeping terms up to \(x^4\):
\begin{align*}
\left(\frac{\sin(x/2)}{x/2}\right)\cos\left(\frac{x}{3}\right)
&=\left(1-\frac{x^2}{24}+\frac{x^4}{1920}\right)\left(1-\frac{x^2}{18}+\frac{x^4}{1944}\right) \\
&=1-\left(\frac{1}{24}+\frac{1}{18}\right)x^2
+\left(\frac{1}{1920}+\frac{1}{1944}+\frac{1}{24}\cdot\frac{1}{18}\right)x^4+\cdots \\
&=1-\frac{7x^2}{72}+\frac{521x^4}{155520}+\cdots
\end{align*}

Hence
\[
\left(\frac{\sin(x/2)}{x/2}\right)\cos\left(\frac{x}{3}\right)
=1-\frac{7x^2}{72}+\frac{521x^4}{155520}+O(x^6)
\]

  \item Using part (a),
\[
\frac{1}{x}\left(\frac{\sin(x/2)}{x/2}\right)\cos\left(\frac{x}{3}\right)
\approx \frac{1}{x}\left(1-\frac{7x^2}{72}+\frac{521x^4}{155520}\right)
= \frac{1}{x}-\frac{7x}{72}+\frac{521x^3}{155520}
\]

Therefore
\begin{align*}
I&=\int_{\pi/6}^{2\pi/3}\left(\frac{1}{x}\left(\frac{\sin(x/2)}{x/2}\right)\cos\left(\frac{x}{3}\right)\right)\,\mathrm{d}x \\
&\approx \int_{\pi/6}^{2\pi/3}\left(\frac{1}{x}-\frac{7x}{72}+\frac{521x^3}{155520}\right)\,\mathrm{d}x \\
&= \left[\ln x-\frac{7x^2}{144}+\frac{521x^4}{622080}\right]_{\pi/6}^{2\pi/3}
\end{align*}

Now simplify:
\begin{align*}
\ln\left(\frac{2\pi}{3}\right)-\ln\left(\frac{\pi}{6}\right)
&=\ln\left(\frac{2\pi/3}{\pi/6}\right) \\
&=\ln 4
\end{align*}

\begin{align*}
-\frac{7}{144}\left[\left(\frac{2\pi}{3}\right)^2-\left(\frac{\pi}{6}\right)^2\right]
&=-\frac{7}{144}\left(\frac{4\pi^2}{9}-\frac{\pi^2}{36}\right) \\
&=-\frac{7}{144}\cdot\frac{15\pi^2}{36} \\
&=-\frac{35\pi^2}{1728}
\end{align*}

\begin{align*}
\frac{521}{622080}\left[\left(\frac{2\pi}{3}\right)^4-\left(\frac{\pi}{6}\right)^4\right]
&=\frac{521}{622080}\left(\frac{16\pi^4}{81}-\frac{\pi^4}{1296}\right) \\
&=\frac{521}{622080}\cdot\frac{255\pi^4}{1296} \\
&=\frac{8857\pi^4}{53747712}
\end{align*}

So
\[
I\approx \ln 4-\frac{35\pi^2}{1728}+\frac{8857\pi^4}{53747712}
\]

Hence
\[
I\approx 1.20244
\]

So the approximation is \(1.20244\) to 5 decimal places.

  \item Using the calculator integration function,
\[
\int_{\pi/6}^{2\pi/3}\left(\frac{1}{x}\left(\frac{\sin(x/2)}{x/2}\right)\cos\left(\frac{x}{3}\right)\right)\,\mathrm{d}x \approx 1.20169
\]

So the calculator value is \(1.20169\).

  \item Comparing the two answers,
\[
1.20244-1.20169=0.00075
\]

So the approximation in part (b) is an overestimate by about \(0.00075\).

It is accurate to \(3\) decimal places, but not to \(4\) decimal places.
\end{enumerate}
\end{tcolorbox}


\newpage

% ---- Question 14 ----
\questionitem
% [Source: _QBT___Solns__Maclaurin_Series, Q14]

\begin{enumerate}[label=\textbf{(\alph*)}]
  \item Find and simplify the first five terms in the Maclaurin series for $e^x$. \marks{2} \\
  \item Hence, or otherwise, write down the first five terms in the Maclaurin series for $e^{-x}$. \marks{1} \\
  \item Use your answers, together with the identity $\cosh x + 1 = 2\cosh^2\left(\frac{x}{2}\right)$, to show that the Maclaurin series for $\cosh^2\left(\frac{x}{2}\right)$ is
  \[
  a + bx^2 + cx^4 + \cdots
  \]
  where $a$, $b$ and $c$ are rational numbers to be determined. \marks{3} \\
\end{enumerate}

\vspace*{10pt}

\begin{tcolorbox}[colback=white, colframe=scarlet, title={\textbf{Solution}}, breakable, grow to left by=6mm, grow to right by=10mm]
\begin{enumerate}[label=\textbf{(\alph*)}]
  \item Using the standard Maclaurin series
\[
e^x=1+x+\frac{x^2}{2!}+\frac{x^3}{3!}+\frac{x^4}{4!}+\cdots
\]
so the first five terms, simplified, are
\[
e^x=1+x+\frac{x^2}{2}+\frac{x^3}{6}+\frac{x^4}{24}+\cdots
\]

  \item Replace $x$ by $-x$ in the series for $e^x$:
\[
e^{-x}=1+(-x)+\frac{(-x)^2}{2}+\frac{(-x)^3}{6}+\frac{(-x)^4}{24}+\cdots
\]
Hence
\[
e^{-x}=1-x+\frac{x^2}{2}-\frac{x^3}{6}+\frac{x^4}{24}+\cdots
\]

  \item Using
\[
\cosh x=\frac{e^x+e^{-x}}{2}
\]
and the expansions from parts (a) and (b),
\begin{align*}
e^x+e^{-x}
&=\left(1+x+\frac{x^2}{2}+\frac{x^3}{6}+\frac{x^4}{24}+\cdots\right)
+\left(1-x+\frac{x^2}{2}-\frac{x^3}{6}+\frac{x^4}{24}+\cdots\right) \\
&=2+x^2+\frac{x^4}{12}+\cdots
\end{align*}
So
\[
\cosh x=\frac{2+x^2+\frac{x^4}{12}+\cdots}{2}
=1+\frac{x^2}{2}+\frac{x^4}{24}+\cdots
\]

Now use the identity
\[
\cosh x+1=2\cosh^2\left(\frac{x}{2}\right)
\]
Substituting the series for $\cosh x$:
\[
2\cosh^2\left(\frac{x}{2}\right)
= \left(1+\frac{x^2}{2}+\frac{x^4}{24}+\cdots\right)+1
=2+\frac{x^2}{2}+\frac{x^4}{24}+\cdots
\]
Divide by $2$:
\[
\cosh^2\left(\frac{x}{2}\right)=1+\frac{x^2}{4}+\frac{x^4}{48}+\cdots
\]

Therefore,
\[
a=1,\qquad b=\frac14,\qquad c=\frac1{48}
\]
\end{enumerate}
\end{tcolorbox}


\newpage

% ---- Question 15 ----
\questionitem
% [Source: _QBT___Solns__Maclaurin_Series, Q15]

\begin{enumerate}[label=\textbf{(\roman*)}]
  \item Use the Maclaurin series for $\ln(1+x)$ and $\ln(1-x)$ to obtain the first three non-zero terms in the Maclaurin series for
  \[
  \ln\left(\frac{(1+x)^2}{1-2x}\right)
  \]
  State the range of validity of this series. \marks{4} \\
  \item Find the value of $x$ in the range of validity for which
  \[
  \frac{(1+x)^2}{1-2x}=2
  \]
  Hence find an approximation to $\ln 2$, giving your answer to three decimal places. \marks{4} \\
\end{enumerate}

\vspace*{10pt}

\begin{tcolorbox}[colback=white, colframe=scarlet, title={\textbf{Solution}}, breakable, grow to left by=6mm, grow to right by=10mm]
\begin{enumerate}[label=\textbf{(\roman*)}]
\item Using
\[
\ln(1+x)=x-\frac{x^2}{2}+\frac{x^3}{3}-\cdots
\]
and
\[
\ln(1-x)=-\left(x+\frac{x^2}{2}+\frac{x^3}{3}+\cdots\right)
\]

we have
\[
\ln\left(\frac{(1+x)^2}{1-2x}\right)=2\ln(1+x)-\ln(1-2x)
\]

Now substitute into each series.

For \(2\ln(1+x)\),
\[
2\ln(1+x)=2\left(x-\frac{x^2}{2}+\frac{x^3}{3}-\cdots\right)
=2x-x^2+\frac{2x^3}{3}-\cdots
\]

For \(-\ln(1-2x)\),
\[
\ln(1-2x)=-\left(2x+\frac{(2x)^2}{2}+\frac{(2x)^3}{3}+\cdots\right)
\]
so
\[
-\ln(1-2x)=2x+\frac{(2x)^2}{2}+\frac{(2x)^3}{3}+\cdots
=2x+2x^2+\frac{8x^3}{3}+\cdots
\]

Hence
\begin{align*}
\ln\left(\frac{(1+x)^2}{1-2x}\right)
&=\left(2x-x^2+\frac{2x^3}{3}\right)+\left(2x+2x^2+\frac{8x^3}{3}\right)+\cdots \\
&=4x+x^2+\frac{10x^3}{3}+\cdots
\end{align*}

So the first three non-zero terms are
\[
4x+x^2+\frac{10x^3}{3}
\]

For validity:

\(\ln(1+x)\) requires \(-1<x\le 1\), and \(\ln(1-2x)\) requires \(-\frac12\le x<\frac12\).

So the combined series is valid for
\[
|x|<\frac12
\]
and it also converges at \(x=-\frac12\).

\item We solve
\[
\frac{(1+x)^2}{1-2x}=2
\]

So
\begin{align*}
(1+x)^2&=2(1-2x) \\
1+2x+x^2&=2-4x \\
x^2+6x-1&=0
\end{align*}

Using the quadratic formula,
\[
x=\frac{-6\pm\sqrt{36+4}}{2}=\frac{-6\pm\sqrt{40}}{2}=-3\pm\sqrt{10}
\]

The root in the range of validity is
\[
x=\sqrt{10}-3
\]
since \(\sqrt{10}-3\approx 0.162\), while \(-3-\sqrt{10}\) is not in the valid range.

Now
\[
\ln 2=\ln\left(\frac{(1+x)^2}{1-2x}\right)
\]
so using the series from part (i),
\[
\ln 2 \approx 4x+x^2+\frac{10x^3}{3}
\]

Substituting \(x=\sqrt{10}-3\),
\[
\ln 2 \approx 4(\sqrt{10}-3)+(\sqrt{10}-3)^2+\frac{10}{3}(\sqrt{10}-3)^3
\]

This simplifies to
\[
\ln 2 \approx \frac{364\sqrt{10}-1149}{3}
\]

Hence
\[
\ln 2 \approx 0.689688\ldots
\]

So, to three decimal places,
\[
\ln 2 \approx 0.690
\]
\end{enumerate}
\end{tcolorbox}


\newpage

% ---- Question 16 ----
\questionitem
% [Source: _QBT___Solns__Maclaurin_Series, Q16]

Find the Maclaurin series for \[e^{e^{x^2}-x}\] up to and including the term in $x^2$. \marks{4}

\vspace*{10pt}

\begin{tcolorbox}[colback=white, colframe=scarlet, title={\textbf{Solution}}, breakable, grow to left by=6mm, grow to right by=10mm]
Using the standard Maclaurin expansion
\[
e^t=1+t+\frac{t^2}{2}+\cdots
\]
we first expand \(e^{x^2}\):

\[
e^{x^2}=1+x^2+\cdots
\]

So the exponent becomes
\[
e^{x^2}-x=1-x+x^2+\cdots
\]

Hence
\[
e^{e^{x^2}-x}=e^{\,1-x+x^2+\cdots}=e\cdot e^{-x+x^2+\cdots}
\]

Let
\[
u=-x+x^2+\cdots
\]

Then
\[
e^u=1+u+\frac{u^2}{2}+\cdots
\]

Now
\[
u^2=(-x+x^2+\cdots)^2=x^2-2x^3+\cdots
\]
so, up to the term in \(x^2\),
\[
u^2=x^2+\cdots
\]

Therefore
\begin{align*}
e^u&=1+\left(-x+x^2\right)+\frac{1}{2}\left(x^2\right)+\cdots\\
&=1-x+\frac{3}{2}x^2+\cdots
\end{align*}

Multiplying by \(e\),
\[
e^{e^{x^2}-x}=e\left(1-x+\frac{3}{2}x^2+\cdots\right)
\]

So the Maclaurin series up to and including the term in \(x^2\) is
\[
e^{e^{x^2}-x}=e-ex+\frac{3e}{2}x^2+\cdots
\]
\end{tcolorbox}


\newpage

% ---- Question 17 ----
\questionitem
% [Source: _QBT___Solns__Maclaurin_Series, Q17]

\begin{enumerate}[label=\textbf{(\alph*)}]
  \item Using the definitions of $\sinh x$ and $\cosh x$, together with the Maclaurin series expansion of $e^x$, find the first three non-zero terms in the Maclaurin series expansion of $x\cosh x-\sinh x$. \marks{3} \\
  \item Hence, by replacing $x$ with $\mathrm{i}x$ in your answer to part (a), find the first three non-zero terms in the Maclaurin series expansion of $x\cos x-\sin x$. \marks{3} \\
\end{enumerate}

\vspace*{10pt}

\begin{tcolorbox}[colback=white, colframe=scarlet, title={\textbf{Solution}}, breakable, grow to left by=6mm, grow to right by=10mm]
\begin{enumerate}[label=\textbf{(\alph*)}]
  \item Using
  \[
  e^x=1+x+\frac{x^2}{2!}+\frac{x^3}{3!}+\frac{x^4}{4!}+\frac{x^5}{5!}+\frac{x^6}{6!}+\frac{x^7}{7!}+\cdots
  \]
  we have
  \[
  e^{-x}=1-x+\frac{x^2}{2!}-\frac{x^3}{3!}+\frac{x^4}{4!}-\frac{x^5}{5!}+\frac{x^6}{6!}-\frac{x^7}{7!}+\cdots
  \]

  From the definitions,
  \[
  \cosh x=\frac{e^x+e^{-x}}{2},\qquad \sinh x=\frac{e^x-e^{-x}}{2}
  \]

  So
  \begin{align*}
  \cosh x
  &=\frac{1}{2}\left[\left(1+x+\frac{x^2}{2!}+\frac{x^3}{3!}+\frac{x^4}{4!}+\cdots\right)
  +\left(1-x+\frac{x^2}{2!}-\frac{x^3}{3!}+\frac{x^4}{4!}+\cdots\right)\right] \\
  &=1+\frac{x^2}{2!}+\frac{x^4}{4!}+\frac{x^6}{6!}+\cdots \\
  &=1+\frac{x^2}{2}+\frac{x^4}{24}+\frac{x^6}{720}+\cdots
  \end{align*}

  and
  \begin{align*}
  \sinh x
  &=\frac{1}{2}\left[\left(1+x+\frac{x^2}{2!}+\frac{x^3}{3!}+\frac{x^4}{4!}+\cdots\right)
  -\left(1-x+\frac{x^2}{2!}-\frac{x^3}{3!}+\frac{x^4}{4!}+\cdots\right)\right] \\
  &=x+\frac{x^3}{3!}+\frac{x^5}{5!}+\frac{x^7}{7!}+\cdots \\
  &=x+\frac{x^3}{6}+\frac{x^5}{120}+\frac{x^7}{5040}+\cdots
  \end{align*}

  Therefore
  \begin{align*}
  x\cosh x-\sinh x
  &=x\left(1+\frac{x^2}{2}+\frac{x^4}{24}+\frac{x^6}{720}+\cdots\right)
  -\left(x+\frac{x^3}{6}+\frac{x^5}{120}+\frac{x^7}{5040}+\cdots\right) \\
  &=\left(x+\frac{x^3}{2}+\frac{x^5}{24}+\frac{x^7}{720}+\cdots\right)
  -\left(x+\frac{x^3}{6}+\frac{x^5}{120}+\frac{x^7}{5040}+\cdots\right) \\
  &=\left(\frac{1}{2}-\frac{1}{6}\right)x^3+\left(\frac{1}{24}-\frac{1}{120}\right)x^5+\left(\frac{1}{720}-\frac{1}{5040}\right)x^7+\cdots \\
  &=\frac{x^3}{3}+\frac{x^5}{30}+\frac{x^7}{840}+\cdots
  \end{align*}

  Hence the first three non-zero terms are
  \[
  x\cosh x-\sinh x=\frac{x^3}{3}+\frac{x^5}{30}+\frac{x^7}{840}+\cdots
  \]

  \item Replace $x$ by $\mathrm{i}x$ in the result from part (a):
  \[
  (\mathrm{i}x)\cosh(\mathrm{i}x)-\sinh(\mathrm{i}x)
  =\frac{(\mathrm{i}x)^3}{3}+\frac{(\mathrm{i}x)^5}{30}+\frac{(\mathrm{i}x)^7}{840}+\cdots
  \]

  Using
  \[
  \cosh(\mathrm{i}x)=\cos x,\qquad \sinh(\mathrm{i}x)=\mathrm{i}\sin x
  \]
  gives
  \[
  \mathrm{i}(x\cos x-\sin x)
  =\frac{(\mathrm{i}x)^3}{3}+\frac{(\mathrm{i}x)^5}{30}+\frac{(\mathrm{i}x)^7}{840}+\cdots
  \]

  Now simplify the powers of $\mathrm{i}$:
  \[
  (\mathrm{i}x)^3=-\mathrm{i}x^3,\qquad
  (\mathrm{i}x)^5=\mathrm{i}x^5,\qquad
  (\mathrm{i}x)^7=-\mathrm{i}x^7
  \]

  So
  \[
  \mathrm{i}(x\cos x-\sin x)
  =-\frac{\mathrm{i}x^3}{3}+\frac{\mathrm{i}x^5}{30}-\frac{\mathrm{i}x^7}{840}+\cdots
  \]

  Dividing through by $\mathrm{i}$,
  \[
  x\cos x-\sin x=-\frac{x^3}{3}+\frac{x^5}{30}-\frac{x^7}{840}+\cdots
  \]

  Hence the first three non-zero terms are
  \[
  x\cos x-\sin x=-\frac{x^3}{3}+\frac{x^5}{30}-\frac{x^7}{840}+\cdots
  \]
\end{enumerate}
\end{tcolorbox}


\end{enumerate}
\end{document}